% Generated draft: requires mathematical and source review.
\documentclass[11pt]{article}
\usepackage[margin=25mm]{geometry}
\usepackage{fontspec}
\setmainfont{Latin Modern Roman}
\usepackage{amsmath,amssymb,mathtools,unicode-math}
\setmathfont{Latin Modern Math}
\usepackage{enumitem,needspace,xurl,hyperref}
\setlength{\parindent}{0pt}
\setlength{\parskip}{5pt}
\setlength{\emergencystretch}{3em}
\newcommand{\N}{\mathbb{N}}
\newcommand{\Z}{\mathbb{Z}}
\newcommand{\Q}{\mathbb{Q}}
\newcommand{\R}{\mathbb{R}}
\newcommand{\C}{\mathbb{C}}
\newcommand{\F}{\mathbb{F}}
\newcommand{\A}{\mathbb{A}}
\newcommand{\PP}{\mathbb{P}}
\newcommand{\E}{\mathbb{E}}
\newcommand{\eps}{\varepsilon}
\newcommand{\dd}{\,\mathrm{d}}
\begin{document}
\section*{UnsolvedMath ID: 30004084}
\textbf{Source problem code:} catalogue entry
\subsection{Definitions and mathematical statement}
The record asks three existence questions concerning injective projective maps. It does not explicitly define ``injected.'' The formulations below use an algebraic morphism that is injective on complex points; they do not require a closed immersion. This interpretation requires confirmation from the source. We use the convention that a variety is integral, so that a smooth projective curve is connected; the record does not explicitly state this convention.

For a complex projective variety \(X\), an injective morphism to \(\PP^m_{\C}\) means a morphism \(f:X\to\PP^m_{\C}\) such that
\[
\forall x,y\in X(\C),\qquad f(x)=f(y)\Longrightarrow x=y.
\]
Under this interpretation, the curve question is:
\[
\text{For every smooth complex projective curve }C,\quad
\text{does there exist an injective morphism }C\longrightarrow\PP^2_{\C}?
\]
The supplied corrected statement also asks the general question:
\[
\text{For every integer }d\geq 1\text{ and every smooth complex projective variety }X
\text{ with }\dim X=d,\quad
\text{does there exist an injective morphism }X\longrightarrow\PP^{2d}_{\C}?
\]
Here the complex ground field is carried over from the curve question; the general question does not explicitly specify its ground field. The displayed formulation concerns positive dimensions; the source does not specify the range of \(d\).

For the Segre--Veronese question, the standard family is defined as follows. Let \(r\geq 1\), let \(n_1,\ldots,n_r\geq 1\) and \(a_1,\ldots,a_r\geq 1\) be integers, and put
\[
D=\sum_{i=1}^{r}n_i,\qquad
N+1=\prod_{i=1}^{r}\binom{n_i+a_i}{a_i}.
\]
The complete linear system of \(\mathcal O(a_1,\ldots,a_r)\) gives the Segre--Veronese embedding
\[
\nu_{\boldsymbol a}:\prod_{i=1}^{r}\PP^{n_i}_{\C}\hookrightarrow\PP^N_{\C};
\qquad S=\operatorname{im}(\nu_{\boldsymbol a}),\quad \dim S=D.
\]
The source does not specify restrictions on these parameters or explicitly state the ground field for this question.

When \(N>2D\), a linear projection to \(\PP^{2D}_{\C}\) is induced by a surjective linear map
\[
q:\C^{N+1}\longrightarrow\C^{2D+1}.
\]
Its center is \(\Lambda=\PP(\ker q)\), of dimension \(N-2D-1\). If \(\Lambda\cap S=\varnothing\), its restriction to \(S\) is a morphism. Thus the question, in the range where a projection to the stated target exists, is whether for every such \(S\) there is a \(q\) with
\[
\PP(\ker q)\cap S=\varnothing,
\qquad
\forall x,y\in S(\C),\quad \PP(q)(x)=\PP(q)(y)\Longrightarrow x=y.
\]
For \(N=2D\), an invertible linear map supplies the target identification. For \(N<2D\), a projection onto \(\PP^{2D}_{\C}\) in this sense does not exist. The source does not explain whether these parameter choices are excluded, whether a linear inclusion into the target is allowed, or whether the intended target has dimension at most \(2D\). This qualification remains unresolved.
\subsection{Short English statement}
The questions ask whether every smooth complex projective curve admits an injective algebraic morphism to the projective plane, whether every smooth projective variety of dimension \(d\) admits one to \(\PP^{2d}\), and whether every Segre--Veronese embedding admits an injective linear projection to a projective space of twice its dimension. The meaning of injection, the ground field in the latter two questions, connectedness conventions, and projection terminology when the ambient dimension is too small require confirmation. These are statement-only questions; the supplied open status is not independently verified.
\subsection{Sources}
\begin{itemize}
\item Supplied local UnsolvedMath registry record, permanent ID \(30004084\), title: \emph{Injective Projective Maps of Curves and Segre--Veronese Varieties}; exact source problem code: \texttt{catalogue entry}. The record belongs to the Oberwolfach Reports Open Problems collection and gives proposed year 2019. No individual proposer is supplied.
\item Supplied report citation: \emph{Singularities and Homological Aspects of Commutative Algebra} (2020), Oberwolfach Reports, pp.~407--464, DOI \texttt{10.4171/owr/2019/6}. \url{https://doi.org/10.4171/owr/2019/6}. The supplied publication year 2020 and the report identifier 6/2019 are retained as given.
\item Supplied literature source: Ballico--Ventura. \url{https://doi.org/10.1007/s11565-020-00343-5}. No title or further bibliographic details are supplied.
\item Supplied literature source: G\"orlach, \emph{Injection dimensions}. \url{https://arxiv.org/abs/1905.11306}.
\item Supplied literature source: Oberwolfach Report 6/2019. \url{https://doi.org/10.4171/owr/2019/6}.
\item Literature assessment dated exactly \texttt{2026-08-22}: the supplied curation labels the problem open and states that the universal injection questions for smooth curves and Segre--Veronese varieties remain open, with literature providing positive families, bounds, and examples. This is supplied dated curation derived from the OWR record, not an independently verified assessment. Its specific assessment does not separately resolve the status of the general smooth-variety question.
\item The registry labels its statement \texttt{corrected\_verified} and reports that the official report was checked and the omitted general smooth-variety question restored. This describes supplied source verification, not verification performed for this draft. The supplied original statement contains only the curve and Segre--Veronese questions; the corrected statement contains all three.
\end{itemize}
\end{document}
